\section{Automazione delle esecuzioni e gestione degli output}

Dopo la validazione funzionale, la campagna sperimentale è stata automatizzata per garantire ripetibilità tra configurazioni e repliche. L’automazione è stata strutturata in tre blocchi: lancio delle simulazioni, raccolta artefatti grezzi, normalizzazione dei dati per il post-processing.

\begin{lstlisting}[caption={Pseudocodice della pipeline di esecuzione},label={lst:pseudocode_pipeline}]
procedure RunCampaign(config_set, n_runs):
    prepare_output_tree()
    for each cfg in config_set do
        for run_id from 0 to n_runs-1 do
            status <- launch_simulation(cfg, run_id)
            if status != OK then
                mark_failure(cfg, run_id)
                stop_campaign()
            end if
            collect_raw_results(cfg, run_id)   # .sca, .vec, logs
        end for
        summarize_config(cfg)
    end for
    export_metadata()  # versioni, seed, timestamp, commit
end procedure
\end{lstlisting}

Un esempio operativo in Bash è riportato di seguito.

\begin{lstlisting}[caption={Template Bash per batch run e raccolta risultati},label={lst:bash_campaign}]
#!/usr/bin/env bash
set -euo pipefail

CONFIGS=("BrakingScenario" "HeterogeneousScenario")
RUNS=10
ROOT_OUT="results/raw"
mkdir -p "$ROOT_OUT"

for cfg in "${CONFIGS[@]}"; do
  for run in $(seq 0 $((RUNS-1))); do
    OUTDIR="$ROOT_OUT/$cfg/run_$run"
    mkdir -p "$OUTDIR"
    echo "[RUN] cfg=$cfg run=$run"
    ./run -u Cmdenv -c "$cfg" -r "$run" > "$OUTDIR/stdout.log" 2>&1
    find . -maxdepth 1 -type f \( -name "*.sca" -o -name "*.vec" \) -exec mv {} "$OUTDIR"/ \;
  done
done
\end{lstlisting}

Per la gestione degli output, oltre ai file numerici sono stati salvati anche metadati di run (configurazione, indice replica, timestamp, versione toolchain). Questa pratica è risultata essenziale per distinguere, nelle analisi successive, differenze dovute alla logica sperimentale da differenze dovute all’ambiente.
